\documentclass[10pt,a4paper]{article}
\usepackage[utf8]{inputenc}
\usepackage[T1]{fontenc}
\usepackage{lmodern}
\usepackage[margin=2.3cm]{geometry}
\usepackage{amsmath,amssymb,booktabs,microtype}
\usepackage[hidelinks]{hyperref}
\title{\textbf{SotSort}: a stable quicksort with a tunable buffer\\[3pt]
\large Design summary and an experimental comparison (version 8-r)}
\author{Lucas Sotomayor}
\date{18 November 2026 --- technical report, draft (revision 4)}
\begin{document}
\maketitle

\begin{abstract}
SotSort is a stable comparison sort whose single code path runs with anything from a fixed
16\,KB buffer to a full $n$-element buffer. Its core is a stable quicksort that partitions blocks
through the buffer and gathers them with rotations. This report summarises the design, states
which parts are borrowed, and compares an implementation against blitsort, fluxsort, glidesort,
driftsort, GrailSort and WikiSort. On one 2-vCPU machine, with 32- and 64-bit integers, the constant-memory mode takes
between 0.19 and 0.75 of blitsort's time on 16 of 22 input classes and between 0.94 and
1.09 on the other six. With a comparator that rules out the vector partition it is
slower than blitsort on 7 of 11 classes, by up to 2.6 times on sorted input. Against
Virtual-Memory Powersort, measured on its authors' own harness with the same auxiliary memory,
SotSort makes about the same number of comparisons on partially sorted input but is 9--24\%
slower on integers at $n=10^7$, and up to 1.6 times slower when moving an element is expensive;
the extra cost is in rotations. All numbers are preliminary.
\end{abstract}

\section{Scope and status}
This report describes and measures one artifact, \texttt{include/sotsort.hpp}: version 8-r
for Sections~\ref{sec:tests} and~\ref{sec:exp}, and version 8-r3, which differs only in how the
routing estimates turn density (Section~\ref{sec:vmps}), for the comparison with
Virtual-Memory Powersort. Both commits are given in Section~\ref{sec:avail}. The proofs (Hoare-style, in Spanish) cover
the algorithms as specified; they are not machine-checked, and the correspondence between the
proofs and this C++ file rests on differential testing (Section~\ref{sec:tests}). I have not
surveyed the literature in depth, so some ideas presented as mine may have prior art.

\section{Design}
\paragraph{Entry.} A linear scan returns on sorted input and reverses strictly descending input.
A second scan counts direction changes; if there are fewer than $n/8$, the input goes to a
natural merge sort, otherwise to the stable quicksort.

\paragraph{Stable partition.} A block of at most $B$ elements is partitioned by writing every
element both to the front of the block and to the buffer and advancing only one index (no
branch); the buffer part is then copied back. For integers with the natural order, a vector
version compares $w$ lanes against the pivot and compresses them (AVX-512 compress, or a
16-row permutation table with AVX2). Partitioned blocks $T_1F_1$ and $T_2F_2$ are fused by
rotating $F_1T_2$, in binary-counter order, so the cut point of each segment travels on a stack
and no binary search is needed.

\paragraph{Quicksort.} The pivot is a copy (a ninther up to $2^{13}$ elements, the median of
about $2n^{1/3}$ samples above), taken at pseudorandom offsets. A lower bound inherited from the
parent detects ``no element is smaller than the pivot''; the elements equal to the pivot are then
separated and are final. Unbalanced steps consume a budget of $\lfloor\lg n\rfloor+1$; when it
runs out, the range is finished by a bottom-up SymMerge sort.

\paragraph{Leaves and merges.} Ranges of at most 128 elements are split into power-of-two chunks
sorted by parity merges between the array and the buffer. A parity merge of unequal lengths $a$
and $b$ needs no bounds checks in its first $\min(a,b)$ iterations. Natural runs are merged with
the power criterion of powersort; each merge is a SymMerge whose leaves use the buffer.

\paragraph{Memory plan.} The fast mode chooses among three tiers from a byte budget: the
built-in buffer, a power-of-two block buffer, or a full buffer of $n$ elements. The budget is
the smaller of a user cap and a fraction of the available memory (including the cgroup margin).
Allocation failure steps down one tier. The same procedures run in every tier with a different
$B$; the proofs only need $B\ge 16$.

\paragraph{Routing heuristics (new in 8-r).} Three cheap probes choose among procedures
that are already proved correct, so they affect time only. (i) Before a buffered merge
on the natural-run path, eight binary searches bound the number of side changes of the
merge; long streaks get a galloping merge, otherwise a parity merge if both runs fit in
the buffer. (ii) When there are few direction changes, 256 samples at pseudorandom
positions tell whether runs have interleaved key ranges; if so, the input goes to the
quicksort unless the number of runs is small. (iii) Eight or more consecutive natural
runs shorter than eight elements are treated as a disordered zone and sorted by the
quicksort as a single run. (ii) and (iii) apply to trivially copyable types only; with
strings they measured worse.

\paragraph{What is borrowed.} The fused triple-reversal rotation (blitsort), the equal-length
parity merge (quadsort), the power criterion (powersort), galloping (Timsort), the lower-bound
treatment of duplicates (pdqsort), the budget fallback (introsort), SymMerge (Kim and Kutzner),
merge path for the parallel mode, and counting sort. \textbf{What I believe is mine}, subject to
the caveat above: the unequal-length parity merge without bounds checks, the stable vector block
partition, the unified constant-to-linear memory design with proofs parametrised by $B$, and the
three routing probes above.

\paragraph{Closest related work.} Moltmann, Nakajima and Wild~\cite{vmpowersort} attack the
same trade-off from the mergesort side: Virtual-Memory Powersort lowers the buffer of an adaptive
stable mergesort from $n/2$ to $O(\sqrt{n\log n})$ elements with the same moves and comparisons
up to an additive $O(n)$. SotSort instead keeps a constant buffer and pays with rotations inside
a quicksort. Section~\ref{sec:vmps} compares the two on their harness. The natural-run path uses the power criterion of
Powersort~\cite{powersort}. Its galloping follows Timsort; Ghasemi, Jugé and
Khalighinejad~\cite{galloping} analyse galloping in this family of natural merge sorts. The
branch-free block partition pursues the goal of BlockQuicksort~\cite{blockqs}, avoiding branch
mispredictions, but keeps stability by writing every element to both destinations.

\section{Testing}\label{sec:tests}
Differential tests against \texttt{std::stable\_sort}, with original indices as tags to check
stability, over 15 input families: 32- and 64-bit signed and unsigned integers, pairs, records of 16, 40 and 128 bytes,
and strings; all three memory tiers; 1 to 4 threads. No failures in 19\,506 cases with AVX-512 and in 7\,410 cases per build with AVX2, no SIMD
and clang, on the final file. AddressSanitizer, UBSan and ThreadSanitizer report nothing.
Allocation failures were injected at every allocation point. Known gap: if the comparator
throws, elements can be lost.

\section{Experiments}\label{sec:exp}
\paragraph{Setup.} One cloud VM with 2 vCPUs (Intel Xeon with AVX-512), 8\,GB, Linux, gcc 13
with \texttt{-O3 -march=native}; Rust 1.97 for glidesort 0.1.2 and the standard library sorts
(\texttt{sort}: driftsort; \texttt{sort\_unstable}: ipnsort). $n=10^6$, median of 5 runs after a
warm-up, pinned to one core. Every output was checked for order, multiset and stability.
Identical runs differed by up to $\pm15\%$ on this machine. The inputs: uniform random; 16 and
$\sqrt n$ distinct values; Zipf ($s=1.1$); sorted; reversed; sorted with 1\% random
replacements; sorted with a 10\% random tail; organ pipe; sorted runs of length $\sqrt n$ and
of length 63 with random keys.

\begin{table}[t]\centering\small
\begin{tabular}{lrrrr}
\toprule
 & \multicolumn{3}{c}{SotSort / blitsort} & SotSort / best of\\
Input & u32 & u64 & key+index & Grail, Wiki (32\,B)\\
\midrule
random & 0.63 & 0.72 & 1.14 & 0.60\\
16 values & 0.46 & 0.62 & 0.80 & 0.41\\
$\sqrt n$ values & 0.19 & 0.49 & 0.86 & 0.48\\
Zipf & 0.44 & 0.50 & 0.64 & 0.50\\
sorted & 1.09 & 1.00 & 2.57 & 0.77\\
reversed & 0.94 & 1.01 & 1.81 & 0.12\\
1\% noise & 1.08 & 1.06 & 1.08 & 0.63\\
10\% random tail & 0.67 & 0.75 & 1.50 & 0.82\\
organ pipe & 0.58 & 0.70 & 1.13 & 0.69\\
runs of $\sqrt n$ & 0.53 & 0.75 & 1.23 & 1.30\\
runs of 63 & 0.42 & 0.48 & 0.78 & 0.75\\
\bottomrule
\end{tabular}
\caption{Constant-memory sorts, $n=10^6$: time ratio, below 1 means SotSort is faster.
``key+index'' is a 64-bit word compared on its upper half (no vector partition). blitsort does
not support 32-byte records.}\label{tab:o1}
\end{table}

\begin{table}[t]\centering\small
\begin{tabular}{lrrrrrr}
\toprule
Input (u64) & SotSort O(1) & SotSort fast & blitsort & fluxsort & glidesort & driftsort\\
\midrule
random & 23.5 & 21.0 & 32.6 & 27.4 & 30.3 & 26.6\\
Zipf & 13.7 & 13.6 & 27.5 & 16.3 & 28.3 & 13.4\\
sorted & 0.3 & 0.3 & 0.3 & 0.3 & 0.6 & 0.4\\
reversed & 0.6 & 0.6 & 0.6 & 0.7 & 0.9 & 0.8\\
1\% noise & 10.8 & 12.1 & 10.2 & 24.2 & 30.5 & 20.1\\
10\% random tail & 5.6 & 7.6 & 7.4 & 6.9 & 5.1 & 6.6\\
organ pipe & 2.8 & 3.6 & 4.0 & 3.2 & 1.9 & 4.9\\
runs of $\sqrt n$ & 25.4 & 22.3 & 33.7 & 26.8 & 12.9 & 44.0\\
runs of 63 & 23.6 & 21.8 & 49.4 & 29.4 & 30.8 & 28.0\\
\bottomrule
\end{tabular}
\caption{Milliseconds, $n=10^6$ 64-bit integers, one thread. ``fast'' uses a full buffer. The
rows with 16 and $\sqrt n$ values are omitted here because the fast mode answers them with a
counting sort.}\label{tab:abs}
\end{table}

\paragraph{Reading.} With 32- and 64-bit integers the constant-memory mode is ahead of
blitsort or level with it on every input class measured; the closest are sorted,
reversed and nearly sorted input (0.94 to 1.09). With a comparator that the vector
partition cannot use (``key+index'') it loses on 7 of 11 classes, worst on sorted input
(2.57), where blitsort's scan is faster. Against the linear-memory sorts, the fast mode
is ahead on random input and behind glidesort on organ pipe, the random tail and runs of
$\sqrt n$ (1.5 to 1.9 times for 64-bit integers). For $n=10^7$ random 64-bit integers
the fast mode took 280\,ms, glidesort 362, the constant-memory mode 364, fluxsort 375,
driftsort 392 and blitsort 463. Extra memory at $n=10^7$ (80\,MB of data), as peak
resident set growth: under 0.5\,MB in constant-memory mode and about 40\,MB in fast mode.

\paragraph{Effect of the routing heuristics.} Against the same code without them
(interleaved runs, minimum of 7): 0.45 to 0.81 of the time on sorted data with a random
zone, 0.26 to 0.88 on sorted runs of random keys, no change on random and
duplicate-heavy input, and at most 10\% worse elsewhere, which is within noise. Buffer
sizes from 4 to 128\,KB and leaf sizes from 32 to 256 made no measurable difference.

\paragraph{Comparisons.} On random input every sort measured makes between 1.07 and 1.29 times
$\lg n!$ comparisons; SotSort makes 1.11. This is not a differentiator.

\paragraph{Adversaries.} McIlroy's adversary, which fixes values lazily, loses to any sort with
a run-detection pass, so a random fraction ($10^{-3}$ to $10^{-1}$) of positions was frozen in
advance. The largest count over $n=10^5,10^6$ was $1.14\,n\lg n$ for SotSort, $0.65$ for
blitsort, $1.96$ for driftsort and $3.67$ for ipnsort. No sort showed quadratic behaviour; this
is evidence, not a proof. A second, lazy adversary that gives every pair of undefined elements
the same minimal value was run against the quicksort alone. With the budget charge on the
equal-separation step removed (the state of version 6), it forces 167, 335, 668 and 1335
comparisons per element for $n=1000$ to $8000$, about $n/6$; with the charge, 20 to 27. That
missing charge was found while writing the termination proof, after fuzzing had not shown it.

\section{Comparison with Virtual-Memory Powersort}\label{sec:vmps}
\paragraph{Setup.} I added SotSort to the experimental harness published with
\cite{vmpowersort} (branch \texttt{assignment-fix}, which also counts copy and move
constructions) and changed nothing else: their inputs (runs of geometrically distributed
length with mean $S$, keys random within a run), their four element types (\texttt{int}, cheap to compare and to move;
\texttt{ptr}, a pointer to 30 integers of which only the last is nonzero, so comparisons are slow
and moves cheap; \texttt{arr}, the same 30 integers held by value, slow in both;
\texttt{fastarr}, 30 random integers by value, cheap comparison and slow moves), their timing with \texttt{clock()}, their correctness
check and \texttt{-O2}. Their wrapper type is not trivially copyable, so SotSort runs its
generic path: no SIMD, no routing probes, no disordered zones. SotSort ran single-threaded in
two budgets: \emph{same memory}, exactly the bytes Virtual-Memory Powersort allocates (plus
SotSort's fixed 16\,KB), and $n/2$ elements, as CPython-style Powersort (\emph{copy}).

\paragraph{A change found by this comparison.} On partially sorted input SotSort first made
about one comparison per element more than Powersort. The cause was the routing step, which
counted every turn before the run detection scanned the array again. Version 8-r3 estimates
the turn density from 64 windows of 256 adjacent pairs when $n\ge 2^{16}$; the estimate only
selects between two correct paths, so an error costs time, not correctness. With it, SotSort's
comparisons are within 2\% of Virtual-Memory Powersort's for $S\ge 100$ at $n=10^7$
(Table~\ref{tab:vm7}) and 9\% fewer for $S=100$ at $n=10^6$; on near-random input ($S=2$)
the quicksort makes 17\% (at $n=10^7$) to 33\% (at $n=10^6$) fewer.

\begin{table}[t]\centering\small
\begin{tabular}{@{}lrrrrrr@{}}\toprule
$S$ & 2 & 100 & 1000 & $10^4$ & $10^5$ & $10^6$\\\midrule
time, SotSort / VM-Powersort, same memory & 1.09 & 1.10 & 1.13 & 1.17 & 1.19 & 1.24\\
time, SotSort / copy-Powersort, $n/2$ memory & 1.06 & 1.10 & 1.09 & 1.05 & 1.10 & 1.05\\
comparisons per element, VM-Powersort & 26.80 & 17.26 & 13.58 & 10.52 & 7.24 & 3.98\\
comparisons per element, SotSort & 22.30 & 17.35 & 13.81 & 10.68 & 7.27 & 4.06\\
moves per element, VM-Powersort & 30.71 & 19.37 & 15.55 & 12.50 & 9.22 & 5.96\\
moves per element, SotSort & 44.99 & 31.96 & 27.04 & 22.51 & 16.89 & 10.77\\\bottomrule
\end{tabular}
\caption{\texttt{int}, $n$ uniform in $[0.9\cdot10^7, 10^7]$ (per-element counts divided by the mean,
$0.95\cdot10^7$), cloud VM, 10 runs for time and 3 for counts. ``Moves'' counts assignments and
copy/move constructions.}\label{tab:vm7}
\end{table}

\paragraph{Results.} With the same memory, Virtual-Memory Powersort is faster on its inputs.
At $n=10^7$ on integers SotSort is 9--24\% slower (Table~\ref{tab:vm7}); with $n/2$ memory it
is 5--10\% slower than copy-Powersort. At $n=10^6$, over the four types, the time ratio to
Virtual-Memory Powersort with the same memory was, on the cloud VM and on a second machine
(a 2-vCPU Linux VM on an Intel Core Ultra 7 155H, gcc 11.4): \texttt{int} 0.87--1.08 and
0.99--1.14; \texttt{ptr} 0.90--1.07 and 0.93--1.14; \texttt{arr} 0.71--1.19 and 0.83--1.11;
\texttt{fastarr} 1.42--1.59 and 1.21--1.39, except $S=10^6$ (0.83 and 0.81). At 10 runs per
cell, ratios moved by up to 7\% between repeated runs, so differences below about 8\% are not
meaningful.

\paragraph{Where the moves go.} Counting moves by phase ($n=950\,000$, \texttt{int}), all the
excess is in rotations inside merges whose two runs both exceed the buffer: 3.75, 3.76 and 4.30
moves per element for $S=1000$, $10^4$ and $10^5$ with the same memory as Virtual-Memory
Powersort, against 0.75, 0.74 and 0.71 with $n/2$; the merges through the buffer cost the same
in both budgets. SymMerge halves the longer run and rotates about half the range per level for
$\log_2(m/B)$ levels, so with a buffer of $O(\sqrt{n\log n})$ elements each large merge pays
$\Theta(\log(m/B))$ extra moves per element, a cost Virtual-Memory Powersort avoids by moving
pages through pointers. Replacing the rotation by a cycle-leader rotation, which moves each
element once, lowered the moves but made 7 of 12 \texttt{fastarr} cells slower, because of its
strided access; it was not kept. Merging larger-than-buffer runs without rotations, for
example by block merging through the buffer, is the open problem this comparison exposes.

\paragraph{Reproduction.} Our absolute times for Virtual-Memory Powersort are close to those
the authors publish for an i7-4770 (e.g.\ 0.129\,s against 0.128\,s for $S=10^6$), but their
data show copy-Powersort 10--25\% faster than Virtual-Memory Powersort, while here the two are
within 4\%. Their algorithms' relative costs are therefore not reproduced on our machine.

\section{Limitations}
For Section~\ref{sec:exp}, one machine, virtualised and shared, with two cores (Section~\ref{sec:vmps}
adds a second virtual machine, still no bare metal); no hardware counters; no ARM; no strings in
the comparison tables; the competitors were built with their default settings; glidesort and
driftsort cross a C-to-Rust call boundary once per sort. The comparison against blitsort uses
its default 512-element buffer, while SotSort's buffer is 16\,KB (2048 64-bit elements): both
are constant, but they are not equal. The thresholds
of the routing probes were tuned on this one machine, on some of the same inputs used for the
comparison, so the structured-input rows are optimistic.

\section{Availability}\label{sec:avail}
Code, tests, the benchmark harness and raw measurements: \url{https://github.com/sotomayorlucas/sotsort},
code at commit \texttt{c16d43f} for version 8-r and \texttt{826125b} for version 8-r3 (MPL-2.0).
The Virtual-Memory Powersort comparison is in \texttt{bench/vmpowersort}: a patch to their
harness, the scripts and the raw results; their code is not redistributed.
\begin{thebibliography}{9}\small
\bibitem{vmpowersort} F.~Moltmann, T.-V.~Nakajima, S.~Wild. Virtual-Memory Powersort. In
  \emph{ESA 2026}, LIPIcs 388, 14:1--14:17, 2026. doi:10.4230/LIPIcs.ESA.2026.14.
\bibitem{powersort} J.~I.~Munro, S.~Wild. Nearly-optimal mergesorts: fast, practical sorting
  methods that optimally adapt to existing runs. In \emph{ESA 2018}, LIPIcs, 2018.
\bibitem{galloping} E.~Ghasemi, V.~Jugé, G.~Khalighinejad. Galloping in fast-growth natural merge
  sorts. In \emph{ICALP 2022}, LIPIcs 229, article 68, 2022.
\bibitem{blockqs} S.~Edelkamp, A.~Weiß. BlockQuicksort: avoiding branch mispredictions in
  Quicksort. In \emph{ESA 2016}, LIPIcs, article 38, 2016.
\end{thebibliography}
\end{document}
